\section{Introduction}\label{sec:intro}

The bulk-realization program of this repository asks for an independently
positive system whose state norm is the localized Weil quadratic form
\cite{InverseBulk}. Every candidate tested so far has been asked to reproduce a
\emph{value}: a weight, a pairing, a norm. That has turned out to be a weak
demand. In the one setting where the search was carried out concretely, the set
of reachable states was everything analytic past the unit circle, so an exhibited
weight identity carries no information unless it comes with a reason the
representation could not have produced a different weight \cite{Bottleneck}.

This paper takes up a proposal that asks for something else. A Wilson line, when
deformed, has an expectation value obeying a differential equation fixed by the
variation of its endpoints together with the non-Abelian Stokes theorem. The
shifted family of the program already obeys a differential equation of exactly
that shape, equation (1.7) of the background \cite{Background}:
\[
\frac{\partial V_{\omega,L}}{\partial\omega}=-G_{\omega,L}V_{\omega,L},
\qquad V_{0,L}=I .
\]
The proposal is that these are the same equation, so that the object to search
for is a \emph{deformation} rather than an operator. The attraction is
structural: a generator and an initial condition determine a flow completely, so
matching a flow leaves nothing to tune. It is the right kind of demand to make.

\subsection*{The shape of what follows}

Pursuing the proposal turns the transfer into a recognizable object, and that
object reorganizes everything else. The arc of the paper is this.

\textbf{First the generator is written out} (Section~\ref{sec:generator}). The
background defined $G_{\omega,L}$ by naming its symbol and inherited the
realization silently from the transfer. Written out, its three factors deliver
the three terms of the target exactly, and the whole shift dependence collapses
to one scalar factor: the shifted form is the central form with every
off-diagonal kernel multiplied by $\cosh(\omega(x-y))$ and every prime atom by
$\cosh(\omega\log n)$, the diagonal unchanged (Theorem~\ref{thm:kernel}). So the
equations of the background carry the \emph{whole} Weil form, not its archimedean
part.

\textbf{Then the transfer is identified} (Section~\ref{sec:allpass}). Its symbol
is \emph{unimodular on the critical line} --- exactly, unconditionally, by the
functional equation and the reality of $\xi$ alone (Proposition~\ref{prop:unimodular}).
The transfer changes no amplitude at any frequency, only phase. It is an
\textbf{all-pass filter}, and everything after this section is a consequence of
that one fact.

Four consequences, in the order they fall out.

\begin{enumerate}
\item \textbf{The phase derivative is the archimedean symbol.} The group delay of
the filter is $2\theta'(\tau)=b(\tau^2)+w_0$
(Proposition~\ref{prop:groupdelay}). The companion investigation found that
identity as a computation between special functions; here it is the statement
that the Weil symbol is a \emph{phase derivative}, which is also why the contact
constant is not a separate ingredient and why a packet at frequency $N$ is held
for a time $\sim\log N$.
\item \textbf{The phase jumps are resonances.} Each zero contributes to
$|K_\omega(i\tau)-1|^2$ a Lorentzian of height $4$ and half-width $\omega$, of
mass $4\pi\omega$, so the Weil form is the total resonant response of the input
(Proposition~\ref{prop:resonance}). The zeros are the filter's resonances.
\item \textbf{Under the Riemann hypothesis the defect is a boundary flux.} The
filter is then inner, the transfer is unitary, and
$\langle f,D_{\omega,L}f\rangle=\int_{L/2}^\infty|(V_\omega f)|^2$
(Proposition~\ref{prop:flux}): Weil positivity on $I_L$ says the flow pushes mass
out of the leading endpoint and never pulls it in.
\item \textbf{A filter that changes only phase is a phase, hence abelian}
(Section~\ref{sec:abelian}). Truncation to an interval is multiplicative on causal
kernels, so the transfers and their generators lie in one commutative algebra and
the path-ordered solution of the flow equation collapses to an ordinary
exponential. \emph{The mechanism the proposal names has, as things stand, nothing
to act on, and any non-commutativity must be supplied transverse to $\omega$.}
\end{enumerate}

\textbf{The shift is then identified geometrically} (Section~\ref{sec:hyperbolic}):
it acts on inputs by the hyperbolic rotation generated by multiplication by $x$,
so it is a deformation and its gauge-theoretic analogue is an angle --- a cusp ---
rather than a coupling that can be dialled.

\textbf{The shift--length plane is then settled completely}
(Section~\ref{sec:region}). The contraction region is monotone in $L$ with
supremum \emph{one or infinity and nothing between}, so contraction at every $L$
is \emph{equivalent} to the absence of zeros at distance more than $\omega$ from
the critical line (Corollary~\ref{cor:graded}). The shifted family is a
\textbf{graded criterion}, one zero-free strip per shift, interpolating to the
Riemann hypothesis as $\omega\downarrow0$. Under that hypothesis the region is the
whole quadrant and has no boundary at all, so there is no critical path to find;
what the plane is good for is the opposite, since monotonicity makes it a
\textbf{disproof instrument} on which only a certified value above one carries
information. The same section locates the much smaller region on which the
\emph{generator} is positive, identifies the quantity that controls its boundary
as an inverse zero-placement precision, and records why that boundary is an
artifact of a sufficient condition rather than a feature of $\zeta$.

\textbf{The margin is then identified, and its collapse traced to the wrong
cause} (Section~\ref{sec:sampling}). Under the Riemann hypothesis
$m_L=\lambda_{\min}(Q_{0,L})$ is the constant in the \emph{lower} sampling
inequality for the Riemann zeros on the Paley--Wiener space $PW_{L/2}$ --- there
is no upper one, because the zeros are not relatively separated
(Proposition~\ref{prop:noupper}) --- and it is strictly positive, because the
divergence of the group delay makes the form domain embed compactly and a counting
argument then forbids a null vector (Proposition~\ref{prop:positive}). So the
unboundedness of the Wigner--Smith delay and the positivity of the margin are one
fact. The zero density meets the Nyquist density at $\tau_c=2\pi e^L$, with an
exact sample deficit $2e^L-\frac74$, so the sampling crossover is the
explicit-formula balance; and the frequency profile of the obstruction is
measurably universal in $\tau/\tau_c$. But the deficit is not what the collapse
measures. Replacing the zeros by the \textbf{Gram points}, which have the same
counting function and no arithmetic in them at all, leaves a margin that collapses
at the same superexponential rate, to within a ratio of logarithms falling from
$2.10$ to $1.05$ across the computed range. \textbf{The collapse is a density
effect, and the margin carries almost no arithmetic information} --- which removes
the reason to ask the mechanism question in the form Section~\ref{sec:region}
poses it, and correspondingly lowers what a candidate source has to reproduce.

\textbf{What remains open is the endpoint direction} (Section~\ref{sec:endpoint}).
There the variation is rank one at the leading end, with the boundary value of
the evolved field as its coefficient; the connection is non-abelian but flat; the
generator's own $L$-variation is its boundary value; and the connection fails to
close on the form domain by exactly the Sobolev borderline. The prime data sits
in the atoms of the connection's kernel, so the surviving object is a flat
connection with punctures at the prime periods.

\textbf{Finally, conditions} (Section~\ref{sec:conditions}). Three of them, one
excluding the displacement operator outright, one satisfied structurally by a
bilinear in unbounded endpoint operators and by no character, and one recording
what the spectrum of dimensions has to carry. Conformal invariance turns out to be
aligned with the target rather than in tension with it, and $L$ is the one
conformal invariant of the configuration.

\subsection*{What this paper does not claim}

No source is constructed. No positivity is proved. No gauge theory is matched,
and none is excluded. Sections~\ref{sec:allpass}--\ref{sec:region} contain
statements conditional on the Riemann hypothesis; each is labelled where it
occurs, and Section~\ref{sec:region} explains structurally why no argument of this
kind could be a route to a proof. Section~\ref{sec:conditions}'s conditions are
proposals, and one of them carries an identification hypothesis the program
deliberately does not make. Proposition~\ref{prop:trace} is a proof sketch resting
on a smoothing estimate quoted from \cite{Background}.

\subsection*{Relation to the companion investigations, and to the earlier program}

The conventions, the target and the criterion are those of
\cite{InverseBulk,Selection}; Section~\ref{sec:conventions} restates them so that
this paper is self-contained apart from the background \cite{Background}, which it
cites by equation number for the shifted family.

Two inherited facts govern the work and are not re-derived. The accumulated
exclusions of this program are statements about \emph{channel decompositions} and
about the reachable states of particular representations; none touches a gauge
theory, and none constrains a flow, which is not a decomposition
\cite{ExclusionMap}. And the contraction whose norm is Weil positivity is
saturated, with a margin collapsing superexponentially in $L$, so no mechanism
may supply a margin uniform in $L$ \cite{Selection}. Section~\ref{sec:region}
shows the flow formulation inherits that constraint rather than escaping it, and
explains the mechanism behind it.

This paper also rejoins an earlier line of the wider program. The shifted-zeta
storage-depth investigation \cite{CriticalPath,CumulativePath} looked for a
critical path in the shift--length plane, developed the hyperbolic identity that
appears here as Proposition~\ref{prop:hyperbolic}, and introduced the Cayley
coordinate in which the contraction criterion becomes positive-realness.
Section~\ref{sec:region} settles the question that investigation was asking and
explains, quantitatively, why its sufficient condition could not reach the region
where anything is detectable.

Finally, the distinction between lines and loops that motivates the whole
investigation is not a matter of open versus closed geometry. The test functions
carry no boundary conditions; the localization is by support alone. What excluded
the closed benchmark was bounded spectrum \cite{Selection,Bottleneck}, and
Section~\ref{sec:conditions} restates that in the form the deformation picture
needs --- as a condition on a group delay.
